\subsection{Why quadratic and projection tests do not settle order four}

The normal-space rank condition does not by itself exclude a six-dimensional
space, even if all forbidden zero rectangles are tested.  Here is an explicit
boundary of that approach.

\begin{proposition}\label{prop:quadratic-projection-boundary}
There is a six-dimensional subspace $L\subseteq\operatorname{Mat}_4(\mathbb C)$
on which $S_1=S_2=0$, such that no matrix $J+t$ with $t\in L$ has an
$r\times s$ zero rectangle with $r+s\geq5$.  The space $L$ is not level.
\end{proposition}

\begin{proof}
Let $\omega$ be a primitive cube root of unity.  Let $H_0$ be the standard
five-dimensional hook: its first row sums to zero, its first column below that
row sums to zero, and all other entries vanish.  Let $B$ have zero first row and
column and lower-right block
\[
 \begin{pmatrix}1\\-1\\0\end{pmatrix}
 \begin{pmatrix}1&\omega&\omega^2\end{pmatrix}.
\]
Put $L=H_0\oplus\mathbb C B$.  The supports make this a direct sum.  The matrix
$B$ is doubly centred and
\[
 \sum_{i,j}B_{ij}^2=(1+1)(1+\omega^2+\omega^4)=0.
\]
Consequently $S_1(B)=S_2(B)=0$.  Its $S_2$ pairing with $H_0$ vanishes because
its row and column sums are zero and the supports are disjoint.  Thus
$S_1=S_2=0$ on $L$.

Write $t=r+c+zB$, with $r$ the first-row part of the hook and $c$ its remaining
first-column part.  In the last three columns, the lower three rows of $J+t$ are
\[
 \begin{pmatrix}
  1+z&1+z\omega&1+z\omega^2\\
  1-z&1-z\omega&1-z\omega^2\\
  1&1&1
 \end{pmatrix}.
\]
The first row sums to four, the three lower entries in the first column sum to
three, and the displayed entries rule out every other zero row or column.
A $2\times3$ rectangle containing the last row has a nonzero entry.  One
containing both middle rows cannot vanish in any nonfirst column.  The remaining
possibility consists of the first row and one middle row; its columns include
two nonfirst columns, whose distinct cube-root coefficients prevent both entries
in that middle row from vanishing.  A $3\times2$ rectangle contains two lower
rows and a nonfirst column, so the same alternatives rule it out.  These are all
minimal rectangles with $r+s=5$, and every larger forbidden rectangle contains
one of them.

Finally take $r=E_{12}-E_{11}$ and $c=E_{21}-E_{41}$.  Then
\[
 t(z)=\begin{pmatrix}
 -1&1&0&0\\
 1&z&z\omega&z\omega^2\\
 0&-z&-z\omega&-z\omega^2\\
 -1&0&0&0
 \end{pmatrix}.
\]
Direct expansion of the matching sums gives
\[
 S_1(t(z))=S_2(t(z))=0,\qquad
 S_3(t(z))=z-2z^2,\qquad S_4(t(z))=2z^2.
\]
For example, the coefficient of $z$ in $S_3$ comes from one entry of each of
$r,c,B$; the coefficient of $z^2$ is the contribution from $r$ and two entries
of $B$.  The only nonzero fourth-order terms use the last row at column one
and the first row at column two, giving $2z^2$.  Since
$\operatorname{per}(J+t)=24+6S_1+2S_2+S_3+S_4$, it follows that
\[
 \operatorname{per}(J+t(z))=24+z.
\]
Thus $L$ is not level.  At $z=-24$ the permanent vanishes despite the absence of
any forbidden zero rectangle.
\end{proof}

The example satisfies the normalized projection-annihilator conditions and the
quadratic Gram-rank condition.  It shows that those conditions detect support
obstructions but do not detect cancellation among permanent terms.  A proof of
the order-four bound must also control the mixed cubic equations.
